\chapter*{阅读指南}
\addcontentsline{toc}{chapter}{阅读指南}
\markboth{阅读指南}{阅读指南}

\section*{编写方式}

本书以操作为主。每章结构相同：给出可直接执行的命令，解释其构成，再提供若干变式练习。理论部分只出现在解释原因的位置，不单独成章。

执行命令时建议手工输入。前两个月的学习效果与输入量相关：复制会跳过拼写错误，也跳过了对命令结构的注意，而拼写错误是理解命令行规则最直接的入口。

\section*{三种安装路线}

读者的硬件条件不同，安装章按三种情况分别给出步骤。

\begin{booktable}{\textwidth}{L{0.782} L{0.854} L{1.493} J{0.871}}
\toprule
路线 & 目标机器 & 覆盖范围 & 章节 \\
\midrule
\endhead
实体机 & 一台可整机安装的计算机 & 安装、显卡、无线、蓝牙、休眠、固件、双系统 & 第\ref{ch:install}章（建议起点） \\
虚拟机 & 现有计算机上的虚拟磁盘 & 安装、配置与回滚全流程，不含硬件 & 第\ref{ch:install-vm}章 \\
VPS & 每月数十元 & 公网服务器、域名、TLS、远程运维 & 第\ref{ch:vps}章 \\
\bottomrule
\end{booktable}

建议顺序为实体机、虚拟机、VPS。实体机给出完整的系统，硬件那一层经验只在实体机上得到；前面几章的实验也以一台可以随时重装的机器为对象。没有可整机安装的计算机时，先用虚拟机把流程走一遍，代价是没有显卡、无线与休眠这些内容。服务器部分放在最后，日常使用熟练之后再到 VPS 上验证。

\begin{warnbox}[实体机先备份]
实体机的分区操作不可撤销。第一次安装前把原有数据复制到外部磁盘，并确认备份可以读取。第\ref{ch:install}章的“确认磁盘与固件”一节给出确认目标盘的命令。
\end{warnbox}

\section*{排版约定}

\begin{itemize}
  \item \texttt{等宽字体}表示命令、文件名与配置项，例如 \texttt{systemctl status sshd}。
  \item 代码块中以 \texttt{\$} 开头的行由普通用户执行；需要管理员权限时写作 \texttt{\$ sudo ...}。
  \item 代码块可直接照抄。块外的说明比分块内的命令更需要阅读。
  \item 带色边框的方框用于练习、注意事项与文化注解，属于正文的一部分。
\end{itemize}

\begin{shellsession}
$ ls /etc/nixos
configuration.nix  hardware-configuration.nix
\end{shellsession}

第一行是输入的命令，第二行是程序的输出。提示符 \texttt{\$} 不属于命令。本书不使用 \texttt{\#} 作为 root 提示符，因为它在代码高亮中会被识别为注释。

\section*{旧工具与新工具}

服务器上始终存在 \texttt{find}、\texttt{grep}、\texttt{sed}；个人工作站可以安装 \texttt{fd}、\texttt{rg}、\texttt{sd}。本书的处理方式是先讲旧工具，再给新工具，理由有两条：任何一台 Linux 上都有旧工具，系统故障时能依赖的只有它们；新工具在日常使用中更短、更快、出错更少。

每个新工具都标注它替代的旧命令。不必一次记住全部，先取两三个常用项，其余在遇到场景时再查。

\begin{toolswap}{grep -r}{rg}
\texttt{rg} 默认递归搜索，默认跳过版本控制忽略的文件，默认输出行号与颜色。完成同样的检索，命令更短，速度通常高一个数量级。
\end{toolswap}

\section*{关于 Unix 传统}

命令行工具的形态由一组四十年的约定决定：单个工具只做一件事，工具之间以文本流连接，成功时静默，配置文件放在固定位置，手册随系统分发。本书在每章开头以一则短注讲一条约定，第\ref{ch:unix-culture}章集中叙述。理解这些约定之后，未见过的命令也能推断出大致用法。

\section*{每章的结构}

章首导语说明该章解决的问题。正文之后是练习，练习不需要提交，也没有标准答案，作用是把读者带回机器前。末尾的检查点列出读完后应具备的能力；其中任何一条达不到，应回到对应小节，而不是继续往下读。

准备一本笔记，记录三类内容：引起困惑的命令、解决问题的步骤、修改过的配置文件。三个月后回看，这份记录比任何教程都更贴合具体问题。
